\textbf{Phần 1. Trắc nghiệm lựa chọn}
\setcounterpageref{ex}{0}

\begin{ex} %[3D1V48]
Làm tròn số $4768$ đến hàng chục ta được số
\choice
{$4760$}
{$4800$}
{$4700$}
{\True $4770$}
\loigiai{
Chữ số hàng đơn vị là 8 ($8 \ge 5$) nên làm tròn lên hàng chục kế tiếp: $4770$.
\textbf{Đáp án của câu hỏi là D}
}
\end{ex}

\begin{ex} %[3D1V48]
Làm tròn số $6325$ đến hàng trăm ta được số
\choice
{$6400$}
{\True $6300$}
{$6330$}
{$6320$}
\loigiai{
Chữ số hàng chục là 2 ($2 < 5$) nên giữ nguyên chữ số hàng trăm: $6300$.
\textbf{Đáp án của câu hỏi là B}
}
\end{ex}

\begin{ex} %[3D1V48]
Một đoạn đường dài $3842\text{ m}$. Bác tài xế nói: 'Đoạn đường dài khoảng $3800\text{ m}$'. Bác tài xế đã làm tròn số đến hàng nào?
\choice
{Hàng chục}
{Hàng nghìn}
{\True Hàng trăm}
{Hàng đơn vị}
\loigiai{
Số $3842$ làm tròn đến hàng trăm được $3800$.
\textbf{Đáp án của câu hỏi là C}
}
\end{ex}

\begin{ex} %[3D1V48]
Số lớn nhất có bốn chữ số mà khi làm tròn đến hàng chục thì được số $5480$ là
\choice
{\True $5484$}
{$5485$}
{$5479$}
{$5489$}
\loigiai{
Các số làm tròn đến hàng chục thành 5480 là từ 5475 đến 5484. Số lớn nhất là $5484$.
\textbf{Đáp án của câu hỏi là A}
}
\end{ex}

\begin{ex} %[3D1V48]
Số bé nhất có bốn chữ số mà khi làm tròn đến hàng trăm thì được số $7500$ là
\choice
{$7400$}
{$7499$}
{$7549$}
{\True $7450$}
\loigiai{
Các số làm tròn đến hàng trăm thành 7500 là từ 7450 đến 7549. Số bé nhất là $7450$.
\textbf{Đáp án của câu hỏi là D}
}
\end{ex}

\begin{ex} %[3D1V48]
Làm tròn số $9985$ đến hàng chục ta được số
\choice
{$9980$}
{\True $9990$}
{$10\,000$}
{$9900$}
\loigiai{
Chữ số hàng đơn vị là 5 ($5 \ge 5$) nên làm tròn lên: $9990$.
\textbf{Đáp án của câu hỏi là B}
}
\end{ex}

\begin{ex} %[3D1V48]
Có bao nhiêu số có bốn chữ số khi làm tròn đến hàng chục thì được số $3460$?
\choice
{$9$ số}
{$5$ số}
{\True $10$ số}
{$8$ số}
\loigiai{
Các số từ 3455 đến 3464 gồm: $3464 - 3455 + 1 = 10$ số.
\textbf{Đáp án của câu hỏi là C}
}
\end{ex}

\begin{ex} %[3D1V48]
Dân số một xã có $8749$ người. Làm tròn đến hàng trăm, dân số xã đó khoảng
\choice
{\True $8700$ người}
{$8800$ người}
{$8750$ người}
{$9000$ người}
\loigiai{
Chữ số hàng chục là 4 ($4 < 5$) nên làm tròn xuống: $8700$ người.
\textbf{Đáp án của câu hỏi là A}
}
\end{ex}

\begin{ex} %[3D1V48]
Khi làm tròn số $1995$ đến hàng trăm ta được số
\choice
{$1900$}
{$1990$}
{$1980$}
{\True $2000$}
\loigiai{
Chữ số hàng chục là 9 ($9 \ge 5$) nên cộng thêm 1 vào hàng trăm: $19 + 1 = 20$ trăm, tức là $2000$.
\textbf{Đáp án của câu hỏi là D}
}
\end{ex}

\begin{ex} %[3D1V48]
Hai số tự nhiên liên tiếp có tổng bằng $4785$. Làm tròn số lớn hơn đến hàng chục ta được số
\choice
{$2400$}
{\True $2390$}
{$2380$}
{$2395$}
\loigiai{
Hai số liên tiếp: số bé $(4785 - 1) : 2 = 2392$, số lớn là $2393$. Làm tròn 2393 đến hàng chục được $2390$.
\textbf{Đáp án của câu hỏi là B}
}
\end{ex}

\begin{ex} %[3D1V48]
Số nào sau đây khi làm tròn đến hàng trăm được $4600$ và khi làm tròn đến hàng chục cũng được số $4600$?
\choice
{$4648$}
{$4550$}
{\True $4598$}
{$4606$}
\loigiai{
Xét số 4598: hàng chục 9 làm tròn đến hàng trăm được 4600; hàng đơn vị 8 làm tròn đến hàng chục thành $4590 + 10 = 4600$.
\textbf{Đáp án của câu hỏi là C}
}
\end{ex}

\begin{ex} %[3D1V48]
Quy tắc làm tròn số: nếu chữ số ở hàng ngay sau hàng làm tròn bé hơn 5 thì ta làm gì?
\choice
{\True Giữ nguyên chữ số của hàng làm tròn và thay các chữ số phía sau bằng chữ số 0}
{Tăng chữ số hàng làm tròn lên 1 đơn vị}
{Giảm chữ số hàng làm tròn đi 1 đơn vị}
{Bỏ hẳn các chữ số phía sau}
\loigiai{
Quy tắc chuẩn khi chữ số liền sau $< 5$.
\textbf{Đáp án của câu hỏi là A}
}
\end{ex}

\textbf{Phần 2. Trắc nghiệm đúng sai}
\setcounterpageref{ex}{0}

\begin{ex} %[3D1V48]
Xét tính đúng sai của quy tắc làm tròn số:
\choiceTF
{\True Khi làm tròn đến hàng chục, ta xét chữ số ở hàng đơn vị.}
{Khi làm tròn đến hàng trăm, ta luôn cộng thêm 1 vào hàng trăm.}
{\True Số 8550 khi làm tròn đến hàng trăm thì được số 8600.}
{Số 1234 làm tròn đến hàng chục được 1240.}
\loigiai{
\begin{itemize}
    \item Ý a) Đúng: Quy tắc làm tròn hàng chục.
    \item Ý b) Sai: Sai vì nếu chữ số hàng chục $< 5$ thì giữ nguyên.
    \item Ý c) Đúng: Vì chữ số hàng chục là 5 nên làm tròn lên.
    \item Ý d) Sai: Sai vì đơn vị là 4 nên làm tròn thành 1230.
\end{itemize}
\begin{itemize}
    \item a) \textbf{Đúng}.
    \item b) \textbf{Sai}.
    \item c) \textbf{Đúng}.
    \item d) \textbf{Sai}.
\end{itemize}
}
\end{ex}

\begin{ex} %[3D1V48]
Cho số $5678$:
\choiceTF
{\True Làm tròn đến hàng chục được số 5680.}
{\True Làm tròn đến hàng trăm được số 5700.}
{Làm tròn đến hàng chục và hàng trăm cho kết quả giống nhau.}
{\True Hiệu giữa số làm tròn đến hàng trăm và số làm tròn đến hàng chục là: $5700 - 5680 = 20$.}
\loigiai{
\begin{itemize}
    \item Ý a) Đúng: Đơn vị là 8 làm tròn lên.
    \item Ý b) Đúng: Hàng chục là 7 làm tròn lên.
    \item Ý c) Sai: Sai vì $5680 \neq 5700$.
    \item Ý d) Đúng: Phép tính đúng.
\end{itemize}
\begin{itemize}
    \item a) \textbf{Đúng}.
    \item b) \textbf{Đúng}.
    \item c) \textbf{Sai}.
    \item d) \textbf{Đúng}.
\end{itemize}
}
\end{ex}

\textbf{Phần 3. Trả lời ngắn}
\setcounterpageref{ex}{0}

\begin{ex} %[3D1V48]
Làm tròn số 7834 đến hàng chục.
\shortans{7830}
\loigiai{
\textbf{Đáp số:} $7830$\\
Chữ số hàng đơn vị là 4 nên làm tròn thành 7830.
}
\end{ex}

\begin{ex} %[3D1V48]
Làm tròn số 8265 đến hàng trăm.
\shortans{8300}
\loigiai{
\textbf{Đáp số:} $8300$\\
Chữ số hàng chục là 6 nên làm tròn lên 8300.
}
\end{ex}

\begin{ex} %[3D1V48]
Tìm số tự nhiên lớn nhất có 4 chữ số khi làm tròn đến hàng trăm được số 3400.
\shortans{3449}
\loigiai{
\textbf{Đáp số:} $3449$\\
Số lớn nhất là 3449.
}
\end{ex}

\begin{ex} %[3D1V48]
Tìm số tự nhiên bé nhất có 4 chữ số khi làm tròn đến hàng chục được số 6520.
\shortans{6515}
\loigiai{
\textbf{Đáp số:} $6515$\\
Số bé nhất là 6515.
}
\end{ex}

\begin{ex} %[3D1V48]
Làm tròn số 4996 đến hàng chục.
\shortans{5000}
\loigiai{
\textbf{Đáp số:} $5000$\\
Hàng đơn vị là 6 làm tròn lên: $4990 + 10 = 5000$.
}
\end{ex}

\begin{ex} %[3D1V48]
Có bao nhiêu số có bốn chữ số khi làm tròn đến hàng trăm thì được số 9000?
\shortans{100}
\loigiai{
\textbf{Đáp số:} $100$\\
Các số từ 8950 đến 9049 có tất cả $9049 - 8950 + 1 = 100$ số.
}
\end{ex}
